\section{Events}
\subsection*{Why Events}

When three coins are tossed but not seen, we still take the sample space as:\\
$\Omega=\{(b_1,b_2,b_3) \mid b_i \in \{H, T\}\}$\\
After uncovering the first (or just after the first toss: is there a difference?)\\
$F=\{\emptyset, \Omega, \{(HHH,(HTH),(HTT),(HHT)\}, \{(T,b_2,b_3): b_i \in \{H,T\}\}\}$\\
After unvovering the second,\\
$\Omega_0=\{H,T\}$\\
$F=\{\emptyset, \Omega, \{H\}\times\{H\}\times\Omega_0\} \text{ plus other stuff and unions of stuff}$
After uncovering the third,
$F=P(\Omega)$

\subsection*{A glimpse of continuum}
Let us model the height of a 'random' person\\
Alice uses a ruler in milimeters.\\
$\Omega=\{\dots,\{1.7745,1.775\},\dots\}$ Event is what my rulers says\\
Bob uses a ruler in micrometers\\
$\Omega=\{\dots,\{1.7777745\},\dots\}$\\
God uses a ruler\\
$\Omega=\R[0,3]$\\
Don't do this set $\Omega=\R$ and let $F = that stuff$

\subsection*{Terminology}
$\Omega$: Sample Space = all possible outcomes\\
$(\Omega,F)$: Measurable space what you can know about what happened\\
$(\Omega, F, \P)$: Probability Space  how likely it is that an event happens\\
Note: If $\Omega$ is finite, one can put the uniform measure, i.e.\\
$\forall \omega \in \Omega, \P(\{w\})=\frac{1}{|\Omega|}$

\subsection*{Properties}
\begin{itemize}
    \item $\P(A^c)=1-\P(A)$
    \item $\P(A)+\P(B)=\P(A\cup B)+\P(A \cap B)$
\end{itemize}

\begin{example}
    Union bound: $\P(\bigcup\limits_n A_n) \leq \sum\limits_n \P(A_n)$\\
    $\bigcup\limits_n A_n = A_1 \cup (A_2\setminus A_1)\cup(A_3\setminus(A_2\cup A_1))\cup\dots$\\
    $\P(\bigcup\limits_n A_n) = \P(A_1)+\P(A_2\setminus A_1)+\P(A_3\setminus(A_2 \cup A_1))+\dots$\\
    Since all the set minus is $\leq \P(A_n)$ thus $\P(\bigcup\limits_n A_n) \leq \sum\limits_n \P(A_n)$
\end{example}

\subsection*{Conditional Probability and Independence}
\begin{definition}
    $A,B\in F$, $\P(B)\neq 0$

    The probability of $A (condition on/given) B$ is $\P(A \mid B)=\dfrac{\P(A\cap B)}{\P(B)}$
\end{definition}

\begin{definition}
    $A_1,A_2,\dots \in F$ and finite subcollection $i_1<i_2<\dots<i_n$

    $\P(A_{i_1}\cap A_{i_2} \cap \dots \cap A_{i_n})=\P(A_{i_1})\P(A_{i_2})\dots\P(A_{i_n})$

    ($A$ and $B$ are independint if $\P(A \cap B)= \P(A)\P(B)$)
\end{definition}

\begin{example}
    \phantom{text}\\
    Independent weather/biased coin (70\% H 30\% T) toss. $\P(rains)=70\%$

    $\Omega = \{ (rains, H), (rains, T), (clear, H), (clear, T)\}$

    (rains) = \{(rains, H), (rains, T)\}

    (heads) = \{(rains, H), (clear, H)\}

    (rains and heads) = (rains) $\cap$ (heads)\\\\
    Soothsayer's count toss: rain if heads, clear if tails

    $\Omega = \{(rains, H), (clear, T)\}$

    $\P(rains)=70\%$\\\\
    Independent and identically distributed (i.i.d) tosses:

    $n$ tosses, $\P(n heads)=0.7^{n}$

    $\P(at least 1 tail)=1-0.7^{n}$
\end{example}

\subsection*{Set Theorem $\longleftrightarrow$ Probability Theory}

$A\cap B \longleftrightarrow A and B$

$A \cup B \longleftrightarrow A or B$

$A^c \longleftrightarrow A does not happen$

$A\setminus \longleftrightarrow A happens B does not happen$

\begin{theorem}[Total Probability Theorem]
    $B_1, B_2, \dots, B_n$ with $\forall \P(B_i)>0$, partitions $\Omega$

    \begin{align*}
        \P(A)&=\P((A\cap B_1)\cup \dots \cup (A \cap B_n))\\
             &=\sum\limits_{i=0}^n \P(A \cap B_i)\\
             &=\sum\limits_{i=0}^n \P(A \mid B_i)\P(B_i)
    \end{align*}
\end{theorem}

\begin{example}
    3 colors of products: 100 blue, 70 red, 30 purple
    
    defective rate: 3\% 9\% 7\%

    partition $\Omega=B\cup R\cup P$

    $\P(B)=\frac{100}{200}, \P(R)=\frac{70}{200}, \P(P)=\frac{30}{200}$

    $\P(D)=\P(D\mid B)\P(B)+\P(D \mid R)\P(R)+\P(D \mid P)\P(P)$

    $\P(B \mid D)=\frac{\P(B \cap D)}{\P(D)}=\frac{\P(D\mid B)\P(B)}{\P(D)}$
\end{example}